\documentclass[10pt,a4paper]{amsart}
\usepackage[no-math]{fontspec}
\usepackage{amsmath,amssymb,mathtools,xfrac,xcolor}
\usepackage[a4paper,margin=22mm]{geometry}
\setmainfont{Latin Modern Roman}[SmallCapsFont=lmromancaps10-regular.otf]
\usepackage[colorlinks=true,linkcolor=blue,citecolor=teal,urlcolor=blue]{hyperref}
\sloppy\newcommand{\ie}{i.e.\ }\newcommand{\resp}{resp.\ }\newcommand{\skpt}{\par\smallskip}
\usepackage[scr=rsfs,cal=euler]{mathalfa}
\multlinegap0pt
\newenvironment{enumeratei}
{\bgroup\def\theenumi{\roman{enumi}}\def\theenumii{\arabic{enumii}}\begin{enumerate}}
{\end{enumerate}\egroup}

\hyphenation{addi-tion-al belong belongs denote exchan-ge exchan-ged exchan-ges}

\newcommand{\psfrac}[2]{\sfrac{(#1)}{#2}}

\newcommand\mto{\mathchoice{\longmapsto}{\mapsto}{\mapsto}{\mapsto}}
\let\oldbigsqcup\bigsqcup
\renewcommand{\bigsqcup}{\mathop{\textstyle\oldbigsqcup}\displaylimits}

\let\epsilon\varepsilon

\DeclareMathOperator{\Ess}{Ess}
\DeclareMathOperator{\Leb}{Leb}

\usepackage{tikz}
\usetikzlibrary{plotmarks}

\theoremstyle{plain}
\newtheorem{theorem}{Theorem}[section]
\newtheorem{lemma}[theorem]{Lemma}
\newtheorem{proposition}[theorem]{Proposition}
\newtheorem{corollary}[theorem]{Corollary}
\newtheorem{claim}{Claim}

\theoremstyle{definition}
\newtheorem{remark}[theorem]{Remark}
\newtheorem{example}[theorem]{Example}

\newcommand{\A}{\mathcal{A}}
\newcommand{\B}{\mathcal{B}}
\newcommand{\I}{\mathcal{I}}

\newcommand{\Z}{\mathbb{Z}}
\newcommand{\N}{\mathbb{N}}
\newcommand{\R}{\mathbb{R}}

\begin{document}
\title[Linear extensions of symmetric IETs]{Ergodicity of centered linear real extensions of every ergodic symmetric interval exchange}
\author{Przemysław Berk\and Frank Trujillo\and Hao Wu}
\date{}\maketitle\setcounter{section}{1}
Let $I$ be a bounded interval, equipped with the Borel $\sigma$-algebra and Lebesgue measure $\lambda_{I}$. Let $T=(\pi,\lambda)$ be an interval exchange transformation given by a permutation $\pi\in S_0^\mathcal A$ and a length vector $\lambda\in\Lambda^{\mathcal A,I}$
{(see Section \ref{sc:IETs} for precise definitions of these objects)}. It~is not difficult to see that $T$ preserves $\lambda_{I}$. We say that a permutation is \emph{symmetric} if and only if for any $i = 1, \dots, d$ $\pi_1\circ\pi_0^{-1}(i)=d+1-i$.

The main objects of study in this article are (real-valued) skew products over interval exchange transformations (IETs). More precisely, if $(X,\mathcal B,\mu)$ is a probability Borel space and $f:X\to \R$ is such that $\int_X f(x)\,d\mu(x)=0$, then a \emph{skew product} $T_f:X\times \R\to X\times \R$ over a measure-preserving map $(X,\mathcal B,\mu,T)$, is a transformation given by
\[
T_f(x,r):=(T(x),x+f(x)).
\]
We will refer to $f$ as a \emph{cocycle}. It~is not difficult to see, that $T_f$ preserves the product measure of $\mu$ on $I$ and the Lebesgue measure $\lambda_{\R}$ on $\R$. We will investigate the ergodic properties of $T_f$ with respect to the measure $\lambda_I\otimes\lambda_{\R}$ when $T$ is either an IET or, more generally, a~product of $n \ge 2$ copies of an IET.

\begin{theorem}\label{thm: ergo}
Let $T$ be an ergodic symmetric IET on $I = [0, 1)$ and let $f(x)=a(x-\sfrac{1}{2})$ for some $a\in\R\setminus\{0\}$. Then, the skew product $T_f:I\times\R\to I\times \R$ is ergodic with respect to $\lambda_I\otimes\lambda_{\R}$.
\end{theorem}
\section{Interval exchange transformations}
\label{sc:IETs}

\subsection{Notations and basic properties} An \emph{interval exchange transformation (IET)} $T$ on a bounded interval $I$ is a piecewise linear bijection of $I$, with finite number of intervals of continuity, on which $T$ acts via translation. {For convenience and without loss of generality, we~assume the interval $I$ to be of the form $[a, b)$, for some $a, b \in \R$, and the IETs to be right-continuous.}

More precisely, there exists $\A$ is an alphabet of $d\in\N$ elements, a~permutation $\pi=\binom{\pi_0}{\pi_1}$ with $\pi_0,\pi_1:\mathcal A\to\{1,\ldots,d\}$ and a collection of left closed and right open subintervals $\{I_{\alpha}\}_{\alpha\in\mathcal A}$ such that and $\bigsqcup_{\alpha\in\A} I_{\alpha}=I$, $T|_{I_{\alpha}}$ acts via translation, and $\pi_0$ and $\pi_1$ describe the order of intervals respectively before and after
action of $T$. It~is easy to check that $T$ preserves Lebesgue measure on $I$ and that the parameters $(\pi,\lambda)$ fully describe the dynamics of $T$, where $\lambda=[\lambda_\alpha]_{\alpha\in\A}:=[|I_{\alpha}|]_{\alpha\in\A}\in \Lambda^{\mathcal A,I}$ is the vector of lengths of intervals $I_{\alpha}$, with $\Lambda^{\mathcal A,I}:=\{\lambda\in\R^{\A}_{+}\mid \sum_{\alpha\in\A}\lambda_{\alpha}=|I|\}$. {Moreover, we~always assume that $\pi$ is \emph{non-reducible}, that is
\[
\pi_1\circ\pi^{-1}_0\{1,\ldots,k\}=\{1,\ldots,k\}\implies k=d.
\]
Otherwise, we~can decompose $T$ into two non-trivial IETs and consider their properties separately. We denote by $S_0^{\A}$ the set of all non-reducible permutations of alphabet~$\A$.
}

For every $\alpha\in\A$, we~denote by $\partial I_{\alpha}$ the left endpoint of $I_{\alpha}$ and by $c_{\alpha}$ its center point. We say that an IET $T$ has a \emph{connection} if there exist $\alpha,\beta\in\A$ with $\pi_0(\beta)\neq 1$, $\pi_1(\alpha)\neq 1$ and $n\in\N_+$ such that
\[
T^{-n}(\partial I_\beta)=\partial I_{\alpha}.
\]
By a connection we often mean the orbit segment $\{T^{-k}(\partial I_\beta)\}_{k=-n,\ldots,0}$.
If such connection exists, we~denote
\[
M(\beta)=M(T,\beta):=\min_{\alpha\in\A}\min\{n\in\N_+\mid T^{-n}(\partial I_\beta)=\partial I_{\alpha}\}.
\]
Otherwise, we~write $M(\beta)=\infty$. Similarly, we~denote
\[
N(\alpha)=N(T,\alpha):=\min_{\beta\in\A}\min\{n\in\N_+\mid T^{n}(\partial I_\alpha)=\partial I_{\beta}\}
\]
and write $N(\alpha)=\infty$ if such connection does not exist. {Note that we always have $T(\partial I_{\pi_1^{-1}(1)})=\partial I_{\pi_0^{-1}(1)}$, a~\emph{trivial} connection. Hence, we~define $M(\pi_0^{-1}(1)):=1$ and $N(\pi_1^{-1}(1)):=1$. {We say that a connection is \emph{irreducible} if it is not a concatenation of two connections, \ie it contains at most two points of the form $\partial I_{\alpha}$, $\alpha\in\A$}.

}

Note that the existence of a non-trivial connection implies that some
non-trivial integer combination of lengths of exchanged intervals is equal to
$0$. Thus, if the length vector is \emph{rationally independent}, i.e., if
$\sum_{\alpha\in\A}r_{\alpha}\lambda_\alpha=0$ for some $(r_{\alpha})_{\alpha
\in \A} \in\mathbb Q^{\A}$ implies that $r_{\alpha}=0$, for all $\alpha\in\A$,
then there cannot be any non-trivial connection. Hence, almost every IET has no non-trivial
connections.

However, in this article we consider the class of \emph{all} ergodic IETs, and, let us recall, the existence of non-trivial connections does not exclude ergodicity. Nevertheless, it~is well-known that if all $\partial I_{\beta}$ are endpoints of connections, then such an IET cannot be ergodic. For the sake of completeness, let us provide a proof of this fact.

\begin{lemma}\label{lem: notallconnections}
Let $T:I\to I$ be an interval exchange transformation given by a permutation $\pi$ and length vector $\lambda$. If for every $\beta\in \A\setminus \{\pi_0^{-1}(1)\}$ we have $M(\beta)<\infty$, then $T$ has only periodic orbits. {More precisely, the base interval $I$ can be decomposed in a finite number of periodic components, given by semi-closed intervals, such that the period is uniform on each of these components.}
\end{lemma}

\begin{proof}
Note that by assumption the set of points
\[
\{T^{n}(\partial I_{\alpha})\mid n\in\Z,\quad\alpha\in \A\}=\{T^{-n}(\partial I_{\alpha})\mid \alpha\in\A\text{ and } 0\le n\le M(\alpha)\}
\]
is finite. Consider the partition given by those points and let $[a,b)$ be an element of this partition.

Note that $T^n$ acts continuously on $[a,b)$ for all $n\in\N$. Indeed, the only possible points of discontinuity of $T$ are $\{\partial I_{\alpha}\}_{\alpha\in\A}$, hence, if for some $n\in\N$ the map $T^n$ did not act continuously on $[a,b)$, there would exist $\beta\in\mathcal A$ such that $\partial I_{\beta}\in T^{n-1}\big([a,b)\big)$, which contradicts the choice of $[a,b)$. In particular, it~follows that $T^n\big([a,b)\big)$ is an interval, for any $n \in \N$.

By Poincaré's recurrence theorem, there exists $N\in\N$ such that
\[
T^{N}\big([a,b)\big)\cap [a,b)\neq \emptyset.
\]
This implies that $T^{N}\big([a,b)\big)=[a,b)$. Indeed, otherwise either $T^{-N}(a)$ or $T^{N}(a)$ belongs to $[a,b)$. Since $a$ is in the orbits of one of the points $\{\partial I_{\alpha}\}_{\alpha\in\A}$, this yields a contradiction.

To sum up, $T^{N}\big([a,b)\big)=[a,b)$ and $T^N$ acts continuously on $[a,b)$. Since $T$ is a piecewise translation, then so is $T^N$. Thus $T^N|_{[a,b)}$ is the identity map on $[a,b)$, which finishes the proof.
\end{proof}

\begin{remark}
By proceeding symmetrically, one can replace in Lemma \ref{lem: notallconnections} the endpoints of connections with their initial points.
\end{remark}

One of the main consequences of the above lemma is the following.

\skpt
\begin{corollary}\label{cor: existenceofinfiniteorbit}
Assume that $T$ is an ergodic IET. Then there exists {$\beta \in\mathcal A \setminus \{\pi_0^{-1}(1)\}$ such that $M(\beta) = +\infty$. }
\end{corollary}

\begin{proof}
Assume, for the sake of contradiction, that $T$ is ergodic but that the conclusion does not hold. Then, by~Lemma \ref{lem: notallconnections}, there exists a non-empty semi-closed interval $[a, b) \subseteq I$ and $N \geq 1$ such that $T^N\mid_{[a, b)}$ is the identity map in $[a, b)$. Therefore, the set $\bigcup_{i = 0}^{N - 1} T^i\big( \big[a, \psfrac{a + b}{2}\big]\big)$ is a non-trivial $T$-invariant set, which contradicts the ergodicity of $T$.
\end{proof}

\subsection{Induced IETs}\label{sc:induced_IETs}
Throughout the proofs of the main results of this paper, we~will often use the first return map of $T$ to a subinterval $J \subseteq I$, {which we denote by $T_J: J \to J$. More precisely, we~define $T_J$ as $x \mto T^{r_J(x)}(x)$, where $r_J: J \to \N$ is given by
\[
 r_J(x):= \min\{n \geq 1 \mid T^n(x) \in J\}.
\]
We sometimes refer to $T_J$ as the \emph{induced map of $T$ to $J$}.

A priori, the map $T_J$ is not necessarily well-defined for \emph{all} points in $J$, although Poincaré's recurrence theorem guarantees that $T_J$ is well-defined in a full Lebesgue measure subset of $J$. However, it~is well known (see, e.g., \cite[\S 3]{veech_interval_1978}) that for any subinterval $J=[a_J,b_J)\subseteq I$ the induced map $T_J$ is an IET of at most $d+2$ intervals, where the possible discontinuities are given by preimages of the discontinuities of $T$ (at most $d - 1$ points) and of the endpoints of $J$ (at most $2$ points, not necessarily disjoint with the previous set).

More precisely, the possible discontinuity points of $T_J$ are given by
\[
\{T^{-m_{J,\alpha}}(\partial I_{\alpha})\}_{\alpha\in \A \setminus \{\pi_0^{-1}(1)\}}, \quad\ m_{J,\alpha}:=\inf\{n\geq 0 \mid T^{-n}(\partial I_{\alpha})\in \mathring{J}\} \text{ \ for } \alpha \in \A,
\]
together with
\[
T^{-m_l}(a_J), \quad m_l :=\inf\{n\ge 0 \mid T^{-n}(a_J)\in \mathring{J}\},
\]
if $a_J$ is different from the left endpoint of $I$, and
\[
T^{-m_r}(b_J), \quad m_r :=\inf\{n\ge 0 \mid T^{-n}(b_J)\in \mathring{J}\},
\]
if $b_J$ is different from the right endpoint of $I$,} where the preimages for
which $m_{J, \alpha}$ (\resp $m_l$ or $m_r$) is $+\infty$ are disregarded.
However, note that if $T$ is minimal, which is the case if $T$ is ergodic (see
Lemma \ref{lem: minimality}), all the above notions are finite.

Moreover, if $T$ is ergodic and $J=[a_J,b_J)\subseteq I$ is chosen so that
\begin{equation}\label{eq: paramofJ}
\begin{gathered}
\text{$a_J=T^{m_0}(\partial I_{\alpha_J})$ and $b_J=T^{n_0}(\partial I_{\beta_J})$, \quad for some $\alpha_J,\beta_J\in\mathcal A$ and $m_0,n_0\in\Z$,} \\
T^{m}(\partial I_{\alpha_J})\notin J, \qquad \text{for any } m \in \{0, \ldots, m_0\} \text{ with } m \neq m_0 \\
T^n(\partial I_{\beta_J})\notin J \qquad \text{for any } n \in \{0, \ldots, n_0\} \text{ with } n \neq n_0,
\end{gathered}
\end{equation}
then the induced map $T_J$ can be seen as an IET of at most $d$ intervals. Indeed, in this case, the discontinuities of $T_J$ belong to the set $
\{T^{-m_{J,\alpha}}(\partial I_{\alpha})\}_{\alpha\in \A\setminus \{\pi_0^{-1}(1)\}}$. Analogously, if $J$ is of the form \eqref{eq: paramofJ} the discontinuities of $T_J^{-1}$ are contained in
\[
\{T^{n_{J,\alpha}}(\partial I_{\alpha})\}_{\alpha\in \A\setminus \{\pi_1^{-1}(1)\}},\!\quad \text{where}\ n_{J,\alpha}:=\min\{n\geq 1\mid T^{n}(\partial I_{\alpha})\in J\} \text{ for any } \alpha \in \A.
\]
If, in addition, $T$ has no non-trivial connections, then the previous two sets have exactly $d - 1$ elements, and $T_J$ can be naturally seen as an IET on $d=\#\A$ intervals and identified with an element $(\pi_J, \lambda_J) \in S_0^\A \times \Lambda^{\A, J}$.

The following simple auxiliary fact tells us what happens if $T$ has non-trivial
connections.

\begin{lemma}\label{lem: alphabetafetrconnections}
Let $T$ be an ergodic interval exchange transformation of $d=\#\A$ intervals and let $J$ be a subinterval of the form \eqref{eq: paramofJ}. Then $T_J$ can be considered as an interval exchange of $d_J$ intervals, where
\[
d_J:= d - \#\{\alpha\in\A\setminus \{\pi_0^{-1}(1)\} \mid m_{J,\alpha}\ge M(\alpha)\}.
\]
In particular, if $J$ does not contain any point from any connection, then $d-d_J$ is equal to the number of non-trivial irreducible connections of $T$.
\end{lemma}

\skpt
\begin{proof}
Since we know that the discontinuities of $T_J$ belong to $\{T^{-m_{J,\alpha}}(\partial I_{\alpha})\}_{\alpha\in \A\setminus \{\pi_0^{-1}(1)\}}$ which has at most $d-1$ elements, to prove that $T_J$ can be seen as an interval exchange of $d_J$ intervals it is sufficient to show that this set has exactly $d_J-1$ elements.

Assume that $\alpha\in\A\setminus \{\pi_0^{-1}(1)\}$ is such that $m_{J,\alpha}\ge M(\alpha)$ and let $\beta\in \A\setminus \{\pi_0^{-1}(1)\}$ be such that $T^{-M(\alpha)}(\partial I_{\alpha})=\partial I_{\beta}$. Then, by~the assumption on $\alpha$, we~have that
\[
T^{-m_{J,\alpha}}(\partial I_{\alpha})=T^{-m_{J,\beta}}(\partial I_{\beta}).
\]
This shows that the connection that ends in the point $\partial I_{\alpha}$ decreases the number of discontinuities of $T_J$ by 1.
To conclude the proof of the first statement it remains to repeat the above reasoning for all $\alpha\in\A$ satisfying $m_{J,\alpha}\ge M(\alpha)$.

To prove the second assertion, first remark that
\[
\#\{\alpha\in\A\setminus \{\pi_0^{-1}(1)\} \mid M(\alpha) < +\infty\}
\]
is equal to the number of non-trivial irreducible connections. Then it suffices to notice that
\[
\#\{\alpha\in\A\setminus \{\pi_0^{-1}(1)\} \mid m_{J,\alpha}\ge M(\alpha)\} = \#\{\alpha\in\A\setminus \{\pi_0^{-1}(1)\} \mid M(\alpha) < +\infty\},
\]
if $J$ does not contain any point from any connection.
\end{proof}

In view of the previous lemma, throughout this work, if $T: I \to I$ is an ergodic IET and $J \subseteq I$ is of the form \eqref{eq: paramofJ}, we~will consider the induced IET $T_J$ as an IET on $d_J$ intervals, where
\begin{align*}
d_J &= 1 + \# \{T^{-m_{J,\alpha}}(\partial I_{\alpha}) \mid \alpha\in \A\setminus \{\pi_0^{-1}(1)\}\\
&= d -\#\{\alpha\in\A\setminus \{\pi_0^{-1}(1)\} \mid m_{J,\alpha}\ge M(\alpha)\} \leq d.
\end{align*}
We will also identify $T_J$ with an element $(\pi_J, \lambda_J) \in S_0^{\A_J} \times \Lambda^{\A_J, J}$ of a possibly smaller alphabet $\A_J$ and denote by $\{I_\gamma^J\}_{\gamma \in \A_J}$ the intervals exchanged by $T_J$.

Let us point out that, in the same way that an IET $T$ with $d$ intervals might have less than $d - 1$ discontinuities, the induced map $T_J$ might have less than $d_J - 1$ discontinuities, that is, some of the points in $
\{T^{-m_{J,\alpha}}(\partial I_{\alpha})\}_{\alpha\in \A\setminus \{\pi_0^{-1}(1)\}}$ might not be real discontinuity points of $T_J$.

Notice that given an ergodic IET $T: I \to I$ and a subinterval $J \subseteq I$ of the form~\eqref{eq: paramofJ}, by~the minimality of $T$ and $T^{-1}$, we~can express $I$ as a disjoint union of the form
\begin{equation}
\label{eq:towers_decomposition}
I = \bigsqcup_{\gamma \in \A_J} \bigsqcup_{i = 0}^{h_\gamma - 1} T^i(I_\gamma^J),
\end{equation}
where, for any $\gamma \in \A_J$, $h_\gamma$ denotes the first return time to $J$ by $T$ of any point in $I_\gamma^J$. We refer to the above decomposition as the \emph{Rokhlin tower decomposition associated with $T_J$ and to $(h_{\gamma})_{\gamma\in \A_J}$ as the corresponding \emph{height} vector}.

We finish this section by recalling a well-known fact, which we prove for completeness.

\begin{lemma}\label{lem: minimality}
If $T: I \to I$ is ergodic with respect to the Lebesgue measure then it is minimal.
\end{lemma}

\begin{proof}
We will show that for every $x,y\in I$ and every $\epsilon>0$ there exists $m\in\N$ such that $|T^{-m}y-x|<2\epsilon$. Take an interval $J:=[x,x+\epsilon)$ and consider the first return map $T_J$. Let $I^J_{\beta}$ be intervals exchanged by $T_J$ and $h_{\beta}$ the corresponding first return times. Since $T$ is ergodic, the set $\tilde I:=\bigcup_{\beta\in \mathcal B}\bigcup_{k=0}^{h_{\beta}-1} T^k(I^J_{\beta})$ is of full Lebesgue measure.

Define $h:=\max\{h_{\beta}\mid\beta\in\mathcal B\}$. Consider the set
\[
C:=\{T^{j}(\partial I_{\alpha})\mid \alpha\in\A,\ j=0,\ldots,h\}.
\]
Pick $0<\delta<\epsilon$ such that $(y,y+\delta]\cap C=\emptyset$. Since $\tilde I$ is of full measure, there exists $\tilde y\in[y,y+\delta]\cap \tilde I$. By the choice of $\delta$ and by the fact that $T$ is right-continuous, the sets $\{T^{-j}[y,\tilde y]\mid j=0,\ldots,h\}$ are a family of pairwise disjoint intervals and $T^{-1}$ acts on each of them by translation. Since $\tilde y\in\tilde I$, there exists $m\le h$ such that $T^{-m}\tilde y\in J$ and $T^{-m}[y,\tilde y]=[T^{m}y,T^{-m}\tilde y]$. Thus
\[
|T^{-m}y,x|<\epsilon+\delta<2\epsilon,
\]
which finishes the proof.
\end{proof}
\subsection{Parametrizing IETs with similar induced maps}
For every IET $T: I \to I$ given by $\pi \in S_0^\A$ and $\lambda\in\Lambda^{\A,I}$, we~consider the set $\Lambda^{\mathcal A,I}_T$ given by
\begin{equation}
\label{eq:nbhd_connections}
\left\{\tilde\lambda\in\Lambda^{\A,I} \left|\;
\begin{array}{l}
M(T_{\tilde{\lambda}},\beta)=M(T,\beta)\ \text{and}\ N(T_{\tilde{\lambda}},\beta)=N(T,\beta), \quad \text{for}\ \beta\in\A, \\[3pt]
T_{\tilde{\lambda}}^{-M(\beta)}(\partial I_\beta^{\tilde\lambda})=\partial I_\gamma^{\tilde\lambda}\iff T^{-M(\beta)}(\partial I_\beta)=\partial I_\gamma,\\
\hspace*{5cm} \text{for $\beta \in \A$ with }M(\beta)<\infty,\\[3pt]
T_{\tilde{\lambda}}^{N(\beta)}(\partial I_\beta^{\tilde\lambda})=\partial I_\gamma^{\tilde\lambda}\iff T^{N(\beta)}(\partial I_\beta)=\partial I_\gamma, \\
\hspace*{5cm}\text{for $\beta \in \A$ with }N(\beta)<\infty.
\end{array}\right\} \right.
\end{equation}

In the above definition, the conditions on $N(\cdot)$ and $M(\cdot)$ are equivalent, nevertheless we write both of them for completeness. The set $\Lambda^{\mathcal A,I}_T$ denotes all length vectors $\tilde \lambda$ in $\Lambda^{\A,I}$ for which the IET $(\pi, \tilde \lambda)$ has the same connection pattern as $T$. Obviously $\lambda\in \Lambda^{\mathcal A,I}_T$. The following proposition is one of the crucial tools used later in the proofs of the main results. In loose words, it~states that by starting with any IET and considering a Rokhlin tower configuration obtained by inducing on a properly chosen interval, we~can obtain a new IET by perturbing the parameters of this configuration, which has the same combinatorial and connection data as the initial map.

\begin{proposition}
\label{prop: unwinding}
Let $T = (\pi, \lambda) \in S_0^\A \times \Lambda^{\A, I}$ be an ergodic IET and let $J \subseteq I$ be a subinterval of the form \eqref{eq: paramofJ} with endpoints $T^{m}(\partial I_{\alpha_J})$, $T^n(\partial I_{\beta_J})$, for some $\alpha_J, \beta_J \in \A$ and $m, n \in \Z$. Assume that $J$ does not contain any point from any connection of $T$ and let
\begin{equation}
\label{eq:initial_towers}
I = \bigsqcup_{\gamma \in \A_J} \bigsqcup_{i = 0}^{h_\gamma - 1} T^i(I_\gamma^J),
\end{equation}
be the associated Rokhlin tower decomposition of $I$.

Then, for any $v \in \R^{\A_J}_+$ satisfying $\sum_{\gamma \in \A_J} v_\gamma h_\gamma = |I|$, there exists $\tilde \lambda \in \Lambda^{\A, I}_T$ such that the IET $\tilde T = (\pi, \tilde \lambda)$ and the interval $\tilde J$ with endpoints $\tilde T^{m}(\partial \tilde I_{\alpha_J})$, $\tilde T^n(\partial \tilde I_{\beta_J})$, where $\{\tilde I_\alpha\}_{\alpha \in \A}$ denote the intervals exchanged by $\tilde T$, satisfy the following.
\begin{itemize}
\item The induced IET $\tilde T_{\tilde J}= (\tilde \pi ^{\tilde J}, \tilde \lambda ^{\tilde J})$ is defined on the alphabet $\A_J$,
\item $\tilde \lambda^{\tilde J} = v$ and $\tilde \pi^{\tilde J} = \pi^J$,
\item {the vectors of heights in the tower decomposition associated with $\tilde T_{\tilde J}$ and $T_J = (\pi^J, \lambda^J)$, as in \eqref{eq:towers_decomposition}, are identical.}
\end{itemize}
\end{proposition}

In the following, given two intervals $J_1, J_2 \subseteq \R$, we~denote
\begin{equation}\label{eq: defofintervalorder}
J_1<J_2\iff x<y \text{ for any }x\in J_1\text{ and any }y\in J_2.
\end{equation}
Notice that given a collection of disjoint intervals, we~can order it according to the relation above.

\begin{proof}[Proof of Proposition \ref{prop: unwinding}]
Fix $v \in \R^{\A_J}_+$ satisfying $\sum_{\gamma \in \A_J} v_\gamma h_\gamma = |I|$. We will define the desired IET $\tilde T$ as follows. First, we~will change the lengths of the intervals in the Rokhlin tower decomposition of $I$ associated with $T_J$, while keeping their order~in~$I$, to express $I$ as a disjoint union of intervals whose lengths are given by $v$. Then we will define a transformation $\tilde T$ on this union so that it defines a Rokhlin tower decomposition for the new transformation. Finally, we~will check that $\tilde T$ Is an IET with the desired properties.

Since $\sum_{\gamma \in \A_J} v_\gamma h_\gamma = |I|$, we~can represent the interval $I$ as the following union:
\begin{equation}
\label{eq:modified_towers}
I = \bigsqcup_{\gamma \in \A_J} \bigsqcup_{i = 0}^{h_\gamma - 1} \tilde I_{\gamma, i}^J,
\end{equation}
where $\tilde I_{\gamma, i}$ is a left-closed right-open interval of length $v_\gamma$ and
\begin{equation}\label{eq:order_preserved}
\tilde I_{i, \gamma}^J < \tilde I_{j, \beta}^J \iff T^i(I_\gamma^J) < T^j(I_\beta^J),
\end{equation}
for every $\gamma,\beta\in \A_J$ and every $i=0,\ldots, h_\gamma -1$, $j=0,\ldots, h_{\beta}-1$.

Let
\[
\tilde J = \bigsqcup_{\gamma \in \A_J} \tilde I_{\gamma, 0}^J.
\]
Notice that since $J = \bigsqcup_{\gamma \in \A_J} I_{\gamma}^J$ is an interval, by~\eqref{eq:order_preserved}, $\tilde J$ is also an interval. Clearly $|J| = \sum_{\gamma \in \A_J} |\tilde I_{\gamma, 0}^J| = \sum_{\gamma \in \A_J} v_\gamma = |v|_1$. Moreover, since $J = \bigsqcup_{\gamma \in \A_J} T^{h_\gamma}(I_{\gamma}^J)$ we can also express $\tilde J$ as a disjoint union of intervals $\{L^J_{\gamma}\}_{\gamma \in \A_J}$ of lengths given by $v$ and such that $\{L^J_{\gamma}\}_{\gamma \in \A_J}$ and $\{T^{h_\gamma}(I_{\gamma}^J)\}_{\gamma \in \A_J}$ have the same order inside $I$.

We define a transformation $\tilde T$ on $I$ by setting, for any $\gamma \in \A$,
\begin{itemize}
\item $\tilde{T}(\tilde I^J_{\gamma,i})=\tilde I^J_{\gamma,i+1}$ for $0\le i<h_{\gamma}-1$, \item $\tilde T(I^J_{\gamma, h_\gamma - 1}) = L_\gamma^J$,
\end{itemize}
and requiring $\tilde T$ to act via a translation when restricted to these subintervals.

Notice that with these definitions the images, by~$T$ and $\tilde T$ respectively, of the intervals in the decompositions \eqref{eq:initial_towers} and \eqref{eq:modified_towers} are ordered in the same way, that is,
\[
 \tilde T(\tilde I_{i, \gamma}^J) < \tilde T(\tilde I_{j, \beta}^J ) \iff T \big( T^i(I_\gamma^J)\big) < T \big( T^j(I_\beta^J) \big),
\]

We will show that $\tilde T$ can be seen as an IET with the same combinatorial data as~$\pi$ and that its length vector belongs to $\Lambda^{\A, I}_T$.

Denote $H:= \sum_{\gamma\in\mathcal \A_J}h_{\gamma}$. Let $\{I_k\}_{k = 0}^{H - 1}$ be the intervals in the tower decomposition~\eqref{eq:initial_towers} ordered according to their order in $I$, i.e.,
\[
I_{k_1}=T^{j_1}\big(I_{\gamma_1}^J\big)\text{ and }I_{k_2}=T^{j_2}\big(I_{\gamma_2}^J\big)\text{ with }k_1<k_2\iff T^{j_1}\big(I_{\gamma_1}^J\big)<T^{j_2}\big(I_{\gamma_2}^J\big),
\]
and let $\{I_k^+\}_{k = 0}^{H - 1}$ be the intervals in the same tower decomposition ordered according to the order of their images, i.e.,
\[
I_{k_1}^+=T^{j_1}\big(I_{\gamma_1}^J\big)\text{ and }I_{k_2}^+=T^{j_2}\big(I_{\gamma_2}^J\big)\text{ with }k_1<k_2\iff T\left(T^{j_1}\big(I_{\gamma_1}^J\big)\right)<T\left(T^{j_2}\big(I_{\gamma_2}^J\big)\right).
\]

Similarly, let $\{\tilde I_k\}_{k = 0}^{H - 1}$ and $\{\tilde I_k^+\}_{k = 0}^{H - 1}$ denote the intervals in the tower decomposition~\eqref{eq:modified_towers} ordered according to their order in $I$ and the order of their images by $\tilde T$ in~$I$, respectively.

For any $\alpha \in \A$ let $0\le k_\alpha< k_{\alpha}' < H$ and $0\le l_\alpha< l_{\alpha}' < H$ be such that
\[
I_{k}\subseteq I_{\alpha}\iff k_{\alpha}\le k\le k_{\alpha}', \qquad I_{\ell}^+\subseteq T(I_{\alpha})\iff \ell_{\alpha}\le \ell\le \ell_{\alpha}'.
\]
Then for every $k_{\alpha}\le k\le k_{\alpha}'$, the IET $T$ acts as a translation on $I_{k}$ by $\sum_{j<\ell_\alpha}|I_j^+|-\sum_{j<k_\alpha}|I_j|$.

We claim that for any $\alpha \in \A$, $\tilde I_{\alpha}:=\bigsqcup_{j=k_\alpha}^{k_\alpha'}\tilde I_j$ is an interval exchanged by $\tilde T$.
Indeed, since the intervals $\tilde I_k$ and $\tilde I_k^+$ are ordered in the same way as the intervals $I_k$ and $I_k^+$, respectively, the transformation $\tilde{T}$ acts on $\tilde I_k$ via translation by $\sum_{j<\ell_\alpha}|\tilde I_j^+|-\sum_{j<k_\alpha}|\tilde I_j|$. Since the translation value does not depend on $k$, $\tilde{T}$ acts as a translation on the whole interval $\tilde I_{\alpha}$.

Therefore, we~can see $\tilde T$ as an IET on $I$ with $\# \A$ intervals. Moreover, $\tilde T$ has the same combinatorics as $T$ since the intervals (and their images) in both tower decompositions are ordered in the same way. Thus we can identify $\tilde T$ with $(\pi, \tilde \lambda)$ for some $\tilde \lambda \in \Lambda^{\A, I}$, and we denote by $\{\tilde I_\alpha\}_{\alpha \in \A}$ the intervals exchanged by $\tilde T$.

Notice that the tower structure associated with $\tilde T_{\tilde J} = (\tilde \pi^{\tilde J}, \tilde \lambda^{\tilde J})$ is given by \eqref{eq:modified_towers}, that~is, if we denote by $\{\tilde I_\gamma^{\tilde J}\}_{\gamma \in \A_J}$ the intervals exchanged by $\tilde T_{\tilde J}$ then $ \tilde I_\gamma^{\tilde J} = \tilde I_{\gamma, 0}$, for every $\gamma \in \A_J$, and we have
\[
 I = \bigsqcup_{\gamma \in \A_J} \bigsqcup_{i = 0}^{h_\gamma - 1} \tilde I_{\gamma, i}^J = \bigsqcup_{\gamma \in \A_J} \bigsqcup_{i = 0}^{h_\gamma - 1} \tilde T(\tilde I_{\gamma}^{\tilde J}).
\]
Since the intervals $\{\tilde T_{\tilde J}(\tilde I_\gamma^{\tilde J})\}_{\gamma \in \A_J} = \{\tilde T^{h_\gamma}(\tilde I_\gamma^{\tilde J})\}_{\gamma \in \A_J} = \{L_\gamma^J\}_{\gamma \in \A}$ have the same order as the intervals $\{T_{J}(I_\gamma^J)\}_{\gamma \in \A_J} = \{T^{h_\gamma}(I_\gamma^J)\}_{\gamma \in \A_J}$ it follows that $\tilde \pi_{\tilde J} = \pi^J$. Finally, since the lengths of the intervals $\{\tilde I_\gamma^{\tilde J}\}_{\gamma \in \A_J}$ are given by $v$, it~follows that $\tilde \lambda^{\tilde J} = v$.

Moreover, notice that the endpoints of the intervals exchanged by $\tilde T$ and those exchanged by $T$ belong to the same tower floors in their respective decompositions, that is,
\[
 \partial \tilde I_\alpha \in \tilde{T}^i(\tilde I_\gamma^{\tilde J}) \iff \tilde I_\alpha \in T^i(I^J_\gamma),
\]
for any $\gamma \in \A_J$ and any $0 \leq i < h_\gamma$. From this, and since $J$ verifies \eqref{eq: paramofJ}, it~follows that~$\tilde J$ is an interval with endpoints $\tilde T^{m}(\partial \tilde I_{\alpha_J})$, $\tilde T^n(\partial \tilde I_{\beta_J})$.

Since $J$ does not contain any point from any connection, then every connection is contained in one of the towers of the form $\bigsqcup_{i=0}^{h_{\gamma}-1}T^i(I_{\gamma}^J)$, for some $\gamma \in \A_J$. Then, it~follows from the definition of $\tilde T$ (and the previous remarks concerning the endpoints of $\tilde T$) that $\tilde T$ possesses the same connection pattern as $T$, that is, $\tilde \lambda \in \Lambda_T^{\A, I}$.
\end{proof}

\begin{remark}
\label{rmk:simplex}
The previous proposition defines a $(d - d' - 1)$-dimensional simplex $\Delta_J \subseteq \Lambda^{\A, I}$ around $\lambda$ such that for any $\tilde \lambda \in \Delta_J$ the IET $(\pi, \tilde\lambda)$ verifies the conclusions of the proposition, where $d'$ denotes the number of non-trivial {irreducible} connections of $(\pi, \lambda)$.

Indeed, It follows from the proof that the map $v \mto \tilde{\lambda}(v)$ given by the previous proposition is linear on the simplex $\tilde \Delta_J:= \big\{v \in \R^{\A_J}_+ \, \big| \, \sum_{\gamma \in \A_J} v_\gamma h_\gamma = |I| \big\} $ and that $\tilde\lambda(\lambda^J) = \lambda$. Moreover, the map is also injective since we can recover $v$ by inducing the IET associated to $(\pi, \tilde\lambda(v))$ to the interval $\tilde J$. Notice that, by~Lemma \ref{lem: alphabetafetrconnections}, the simplex~$\tilde \Delta_J$ has dimension $d - d' -1$. Thus $\Delta_J = \tilde\lambda(\tilde \Delta_J)$ satisfies the statement above.
\end{remark}

\section{Symmetric interval exchange transformations}

\subsection{Notations and basic properties}
Let $I = [a, b)$ be a bounded interval and $T=(\pi, \lambda) \in S_0^\A \times \Lambda^{\A, I}$ be an IET on $I$ with $d := \#\A$ intervals. The permutation $\pi$ (and any IET having $\pi$ as permutation) is said to be \textit{symmetric} if $\pi_{1}\circ \pi^{-1}_{0}(i)=d+1-i$, for any $1 \leq i \leq d$. We say that $T$ is \emph{non-degenerate} if $\partial I_{\alpha}$ is a discontinuity of $T$, for every $\alpha\in \A {\setminus \{\pi_0^{-1}(1)\}}$. If $T$ is \emph{degenerate}, i.e., if there exists $\partial I_{\alpha}$ which is not a real discontinuity, then whenever we refer to \emph{intervals of continuity} of $T$, we~mean maximal intervals of continuity. Notice that the inverse of a symmetric IET is also a symmetric IET.

We denote the \emph{symmetric reflection} or \emph{involution} on the open interval $\mathring{I} = (a, b)$ by $\mathcal{I}_I$, where $\mathcal{I}_I: \mathring{I}\to \mathring{I}$ is given by $\mathcal{I}_I(x)= a + b - x$. We omit the endpoints of the interval in this definition so that the domain and codomain of the involution are well-defined subsets of $I$.
It is well-known and easy to verify that if $T$ is a symmetric IET then
\begin{equation}\label{eq: standardconjugacy}
\mathcal{I}_I\circ T(x)=T^{-1}\circ\mathcal{I}_I(x) \quad \text{if } x \neq \partial I_{\alpha} \text{ for any } \alpha \in \A.
\end{equation}

\pagebreak[2]
More generally, the equation above implies
\begin{multline}
\label{eq:generalized_conjugacy}
\mathcal{I}_I\circ T^n(x)=T^{-n}\circ\mathcal{I}_I(x) \quad \text{if } x \neq T^{-i}(\partial I_\alpha) \text{ for any } \alpha \in \A \text{ and}\\ \text{any $i$ such that } 0 \leq i < n.
\end{multline}
{Notice that $\I_I\circ T$ and $T^{-1}\circ \I_I$ are not defined everywhere on $I$ since the $\I_I\circ T$ is not defined at $\partial I_{\pi_0^{-1}(d)}$ while $T^{-1}\circ \I_I$ is not defined at $\partial I_{\pi_0^{-1}(1)}$.}
Moreover, {a direct calculation shows that}
\begin{equation}
\label{eq:sym_identity_endpoints}
\I_I\circ T(\partial I_{\alpha})=\partial I_{\hat\alpha},\ \text{where}\ \pi_0(\hat\alpha)=\pi_0(\alpha)+1,
\end{equation}
for $\alpha\in\A$ with $\pi_0(\alpha)\neq d$, and
\begin{equation}\label{eq:sym_identity_endpoints_2}
T^{-1}\circ \I_I(\partial I_{\alpha})=\partial I_{\hat\alpha},\ \text{where}\ \pi_0(\hat\alpha)=\pi_0(\alpha)-1,
\end{equation}
for $\alpha\in\A$ with $\pi_0(\alpha)\neq 1$.

Let us point out that $\pi$ being symmetric is not a necessary condition for \eqref{eq: standardconjugacy} to hold. Indeed, if $\pi$ is not symmetric but its intervals of continuity are exchanged symmetrically {(e.g., by~adding a `fake discontinuity' in one of the exchanged intervals of a symmetric IET with $d$ intervals and considering it as an IET on $d + 1$ intervals)} then \eqref{eq: standardconjugacy} still holds. Moreover, there are examples of IETs that do not exchange their intervals of continuity symmetrically but still satisfy \eqref{eq: standardconjugacy}. These examples arise from two-covers of quadratic differentials {(see, e.g., \cite{boissy_dynamics_2009}}).

The following lemma provides a simple sufficient condition for an IET satisfying~\eqref{eq: standardconjugacy} to be symmetric.

\begin{lemma}
\label{lem: condition for symmetric iets,}
Let $T=(\pi,\lambda): I \to I$ be a non-degenerate IET. If $T$ satisfies
\eqref{eq: standardconjugacy}
and $\lambda_{\alpha}\neq \lambda_{\beta}$ for any distinct $\alpha, \beta \in \mathcal{A}$, then $T$ is symmetric.
\end{lemma}

\begin{proof}
Arrange the intervals $\{I_{\alpha}\}_{\alpha\in \A}$ according to their lengths, such that $\lambda_{\alpha_1}>\lambda_{\alpha_2}>\dots> \lambda_{\alpha_d}$.
Note that $\I_I\circ T$ is continuous on~$\mathring I_{\alpha_1}$, so $T^{-1}\circ \I$ is also continuous on~$\mathring I_{\alpha_1}$. Since the discontinuity points of $T^{-1}$ are given by $\{T(\partial I_{\beta}) \mid \beta\in \A {\setminus \{\pi_1^{-1}(1)\}}\}$, so by the maximal length of $I_{\alpha_1}$, non-degenericity of $T$ and the continuity of $T^{-1}\circ \I_I$, we~must have
\begin{equation}\label{symmetric identity}
\I_I(I_{\alpha_1})=T(I_{\alpha_1}).
\end{equation}

By induction on the index $\{\alpha_i, 1\le i\le d\}$, the above identity holds for every $\alpha \in \A$. Because the involution $\I$ reverses the order of $\{I_{\alpha}, \alpha \in \A\}$, that is:
\[
I_{\alpha}<I_{\beta}\implies \I_I(I_{\alpha})>\I_I(I_{\beta}),
\]
{where the order of intervals is given by \eqref{eq: defofintervalorder}, }the identity \eqref{symmetric identity} implies that $T$ also reverses the order of $\{I_{\alpha}, \alpha \in \A\}$, hence $T$ is symmetric.
\end{proof}
We have an immediate consequence for IETs with general combinatorial data.
\begin{corollary}\label{cor: symmetricintcont}
Let $T: I \to I$ be an IET satisfying \eqref{eq: standardconjugacy} such that all its continuity intervals are of different lengths. Then $T$ exchanges its continuity intervals symmetrically.
\end{corollary}

\begin{proof}
The result follows from Lemma \ref{lem: condition for symmetric iets,} by replacing the intervals exchanged by~$T$ with the continuity intervals of $T$ (thus possibly reducing the number of exchanged intervals).
\end{proof}

Given an IET $T: I \to I$ with exchanged intervals $\{I_{\alpha}\}_{\alpha \in \A}$ we denote by $c_{\alpha}$ the middle point of the interval $I_\alpha$, for any $\alpha\in\A$. Note that if $T$ is symmetric, then\vspace*{-3pt}
\begin{equation}\label{eq: centersandinvolution}
\mathcal{I}_I\circ T(c_{\alpha})=c_{\alpha},
\end{equation}
for every $\alpha\in\A$. These points, as well as the middle point of the interval $I$, which we denote by\vspace*{-3pt}
\[
c_{1/2} := \frac{a + b}{2},
\]
will play an important role in our proofs. Notice that $c_{1/2}$ is the only fixed point of $\I_I$ while the points $\{c_{\alpha}\}_{\alpha\in\A}$ are the only ones satisfying
\eqref{eq: centersandinvolution}. Moreover, the backward and forward
iterates of these points are closely related.

\begin{lemma}
\label{lem:inverse_iterates}
Let $T: I \to I$ be a symmetric IET and $\alpha \in \A \cup \{\sfrac{1}{2}\}$. If $m \geq 1$ is such that $\{c_\alpha, T(c_\alpha), \ldots, T^{m - 1}(c_\alpha)\} \cap \{\partial I_\beta\}_{\beta \in \A} = \emptyset$, then\vspace*{-3pt}
\begin{equation}
\label{eq:inverse_iterates}
T^{-m}(c_\alpha) =
{T^{-1} \circ \I_I} \big( T^{m - \delta_{1/2}(\alpha)}(c_\alpha) \big),
\end{equation}
where\vspace*{-3pt}
\begin{equation}
\label{eq:delta_1/2}
\delta_{1/2}(\beta) = \left\{\begin{array}{ll} 1 & \text{ if } \quad \beta = \sfrac{1}{2}, \\
0 & \text{ otherwise}.
\end{array} \right.
\end{equation}
\end{lemma}

\begin{proof}
This result follows by directly applying \eqref{eq:generalized_conjugacy} and noticing that $\I_I(c_{1/2}) = c_{1/2}$ and that $\I_I(c_\alpha) = T(c_\alpha)$, for any $\alpha \in \A$.
\end{proof}

The following result concerning Birkhoff sums over symmetric IETs follows
directly from \cite[Lem.\,3.11]{berk_ergodicity_2023}. Although not
immediately useful for us, this fact will be crucial in one of the final
arguments of the proof of the main result of this paper.
\begin{lemma}\label{lem: berktrujillo}
Let $T: I \to I$ be a symmetric IET, $\alpha \in \A \cup\{\sfrac{1}{2}\}$ and $f : I\to \R$ satisfying $f\circ \I_I=-f$ on $\mathring I$. If $c_{\alpha}$ is not part of any connection, then\vspace*{-3pt}
\[
S_{2n + \delta_{1/2}(\alpha)}f(T^{-n}(c_{\alpha})) = 0
\]
for any $n\in\N$, where $\delta_{1/2}$ is given by \eqref{eq:delta_1/2}.
\end{lemma}

\subsection{Induced symmetric IETs}
Given an IET $T: I \to I$ and a subinterval $J \subseteq I$ with associated induced map $T_J = (\pi^J, \lambda^J) \in S_0^{\A_J} \times \Lambda^{\A_J, J}$ (see Section \ref{sc:induced_IETs} and Lemma \ref{lem: alphabetafetrconnections}), we~denote by $\{I^J_\gamma\}_{\gamma \in \A_J}$ the intervals exchanged intervals by $T_J$ and by $\{c^J_\gamma\}_{\gamma \in \A_J}$ the middle points of these intervals. We denote by $p_J : I \to J$ the first return map to $J$ by $T^{-1}$, that is, $p_J$ is given by $x \mto T^{-b_J(x)}(x)$, where
\begin{equation}
\label{eq:backward_return}
b_J(x) := \min\{m \geq 1 \mid T^{-m}(x) \in J\}.
\end{equation}
We say that a subinterval $J \subseteq I$ is \emph{symmetric} if there exists $\alpha \in \A$ and $\Delta > 0$ such that $J = [c_\alpha - \Delta, c_\alpha + \Delta) \subseteq I_\alpha$ or $J=\big[c_{1/2}-\Delta, c_{1/2} +\Delta\big)\subseteq I$. To differentiate between the two cases we will refer to the former as \emph{$\alpha$-symmetric} and to the latter as \emph{$\sfrac{1}{2}$-symmetric}.

As we shall see, inducing a symmetric IET on symmetric intervals defines IETs satisfying \eqref{eq: standardconjugacy}, which, as seen in the previous section, is closely related with the symmetry of IETs (see Lemma \ref{lem:symmetric_condition}). Moreover, it~is possible to construct $\alpha$-symmetric subintervals of the form \eqref{eq: paramofJ}, for any $\alpha \in \A \cup \{\sfrac{1}{2}\}$ (see Lemma \ref{lem: sym_endpoints}). Under additional conditions on the symmetric subinterval, we~can guarantee that associated induced maps are actually symmetric.

\begin{proposition}
\label{prop:induced_symmetric}
Let $T: I \to I$ be an ergodic symmetric IET and fix $\alpha \in \A \setminus \{\pi_0^{-1}(1)\}$ such that $M(\alpha) = +\infty$ (see Corollary \ref{cor: existenceofinfiniteorbit}). Fix $m \geq 1$ and let $J \subseteq I$ be the left-closed right-open subinterval with endpoints $T^{-m}(\partial I_\alpha)$, $T^{m}(\partial I_{\hat \alpha})$ (\resp $T^{-m + 1}(\partial I_\alpha)$, $T^{m}(\partial I_{\hat \alpha})$), where $\pi_0(\hat \alpha) = \pi_0(\alpha) - 1$.
Then $J$ is $\beta$-symmetric for some $\beta \in \A$ (\resp $\sfrac{1}{2}$-symmetric) and $T_J$ satisfies \eqref{eq: standardconjugacy}.

If, in addition, $J$ does not contain points from any connection. Then $T_J$ is an ergodic symmetric IET on $d - d'$ intervals, where $d'$ is the number of non-trivial {irreducible} connections of $T$.
\end{proposition}

In the setting of the previous proposition, the exchanged intervals' middle points of the induced map and of the original IET will be closely related.

\begin{proposition}
\label{prop:induced_centers}
Let $T$, $\alpha, \hat \alpha, \beta$ and $J$ as in Proposition \ref{prop:induced_symmetric}. Assume that $J$ does not contain points from any connection. Then the following holds:
\begin{enumerate}
\item For any $\gamma \in \A_J$ there exists unique $\sigma \in \A \cup \{\sfrac{1}{2}\}$, with $ \sigma \neq \beta$, and $\ell \geq 1$ such that $c_\gamma^J = p_J(c_\sigma) = T^{-\ell}(c_\sigma)$ and $T_J(p_J(c_\sigma)) = T^{\ell - \delta_{1/2}(\sigma)}(c_\sigma)$, where $\delta_{1/2}$ is given by \eqref{eq:delta_1/2}. Moreover, $c_\sigma$ does not belong to any non-trivial connection.
\item For any $\delta \in \A \cup \{\sfrac{1}{2}\}$ with $\delta \neq \beta$, either $p_J(c_\delta) = c_\gamma^J$ or $p_J(c_\delta) = \partial I_\gamma^J$ for some $\gamma \in \A_J$. In the latter case, $c_\delta$ lies inside a connection of $T$.
\end{enumerate}
\end{proposition}

For the sake of clarity, we~postpone the proofs of Propositions \ref{prop:induced_symmetric} and \ref{prop:induced_centers} to the end of this section and rather start by proving several preliminary lemmas.

Given a symmetric IET $T: I \to I$, it~is easy to verify that the interior of $\sfrac{1}{2}$\nobreak-sym\-metric intervals are invariant by $\I_I$. On the other hand, as we shall see below, the interior of $\alpha$-symmetric intervals are invariant by $\I_I \circ T$.

\begin{lemma}\label{lem: local_reflection}
Let $T : I \to I$ be a symmetric IET and $J \subseteq I$ be an $\alpha$-symmetric interval, for some $\alpha \in \A$. Then $\mathcal{I}_I\circ T(\mathring{J}) = \mathring{J}$. Moreover,
\[
\I_J = \mathcal{I}_I\circ T|_{\mathring{J}},
\]
that is, $\mathcal{I}_I\circ T|_{\mathring{J}}$ is the symmetric reflection on $J$.
\end{lemma}

\begin{proof}
Let $\alpha \in \A$ and $\Delta > 0$ such that $J = [c_\alpha - \Delta, c_\alpha + \Delta) \subseteq I_\alpha$. Fix $x\in J$ and let $-\Delta < \delta < \Delta$ be such that $x = c_{\alpha}+\delta$. Notice that $\I_J(x) = c_\alpha - \delta$.

Since $T|_{I_{\alpha}}$ is a translation, by~\eqref{eq: centersandinvolution}, we~have
\[
\I_I\circ T(x)=\I_I\circ T(c_{\alpha}+\delta)=\I_I\left(T(c_{\alpha}) +\delta\right)=
\I_I\circ T(c_{\alpha})-\delta=c_{\alpha}-\delta\,\in\, J,
\]
which finishes the proof.
\end{proof}

Notice that any of the exchanged intervals of a symmetric IET $T: I \to I$ defines a symmetric interval. Hence, we~obtain the following as a direct consequence of Lem\-mas~\ref{lem:inverse_iterates} and \ref{lem: local_reflection}.

\begin{corollary}
Let $T : I \to I$ be a symmetric IET, $\alpha \in \A$ and $m \geq 1$. Then\vspace*{-3pt}
\[
\{c_\alpha, \ldots, T^{m - 1}(c_\alpha)\} \cap \{\partial I_\beta\}_{\beta \in \A} = \emptyset\iff \{c_\alpha, \ldots, T^{-m + 1}(c_\alpha)\} \cap \{\partial I_\beta\}_{\beta \in \A} = \emptyset.
\]
\end{corollary}

We now relate the induced IET on a symmetric interval $J$ with the associated involution $\I_J$.

\begin{lemma}
\label{lem:symmetric_condition}
Let $T : I \to I$ be a symmetric IET and $J \subseteq I$ be a symmetric interval. Then $T_J$ satisfies \eqref{eq: standardconjugacy}.\end{lemma}

\begin{proof}
Assume first that $J$ is $\alpha$-symmetric for some $\alpha\in\A$. Let $x\in J$ be not an endpoint of an interval exchanged by $T_J$ and let $h\in\N$ be such that $T_J(x)=T^{h}(x)$. Notice that since $x$ is not an endpoint of the exchanged intervals, by~\eqref{eq:generalized_conjugacy} applied to~$T$,\vspace*{-3pt}
\[
{(\I_I\circ T)\circ T_{J}(x)} = (\I_I\circ T) \circ T^h(x)=T^{-h}\circ (\I_I\circ T)(x).
\]
Thus, to prove \eqref{eq: standardconjugacy}, it~suffices to notice that\vspace*{-3pt}
\[
T_{J}^{-1}\circ (\I_I\circ T)(x) = T^{-h}\circ (\I_I\circ T)(x).
\]
Indeed, by~Lemma \ref{lem: local_reflection}, for any $n \geq 1$, $T^n(x) \in J$ if and only if $(\I_I\circ T) \circ T^n(x) \in J$. Thus the equation above follows since $T_J(x)=T^{h}(x)$ and $T^{-h}\circ (\I_I\circ T)(x) = (\I_I\circ T) \circ T^h(x)$.

Assume now that $J$ is $\sfrac{1}{2}$-symmetric. Let $x\in J$ be not an endpoint of an interval exchanged by $T_J$ and let $h\in\N$ be such that $T_J(x)=T^{h}(x)$. Again by \eqref{eq:generalized_conjugacy} we get\vspace*{-3pt}
\[
\I_I\circ T_{J}(x) = \I_I\circ T^h(x)=T^{-h}\circ \I_I(x)
\]
and, similarly to the previous case, to show \eqref{eq: standardconjugacy} it is sufficient to notice that $T_{J}^{-1}\circ\nobreak\I_I(x)=T^{-h}\circ\I_I(x)$.
\end{proof}

By the result above and given Corollary \ref{cor: symmetricintcont}, if $J$ is symmetric and the lengths of the continuity intervals of $T_J$ are pairwise distinct, then $T_J$ exchanges its continuity intervals symmetrically.

The following lemma shows that if the orbit of a middle point intersects the discontinuities of the IET, then the middle point must be part of a connection.

\begin{lemma}
\label{lem:symmetric_connection}
Let $T: I \to I$ be a symmetric IET and $\beta \in \A \cup \{\sfrac{1}{2}\}$. Assume there exists $m \in \Z$ and $\alpha \in \A$ such that $T^m(c_\beta) = \partial I_\alpha$ and $T^k(c_\beta) \notin \{\partial I_\gamma\}_{\gamma \in \A}$, for any $-|m| < k < |m|$.

Then {if $m\ge 0$}, we~have $\alpha \neq \pi_0^{-1}(1)$ and\vspace*{-3pt}
\[
T^{-m - \delta_{1/2}(\beta)}(c_\beta) = \partial I_{\hat \alpha},
\]
where $\delta_{1/2}$ is given by \eqref{eq:delta_1/2} and $\pi_0( \hat \alpha) = \pi_0(\alpha) - 1$.

{If on the other hand $m<0$, then $\alpha\neq \pi_1^{-1}(1)$ and\vspace*{-3pt}
\[
T^{-m - \delta_{1/2}(\beta)}(c_\beta) = \partial I_{\hat \alpha},
\]
with $\pi_0( \hat \alpha) = \pi_0(\alpha) + 1$.
}

In particular, $c_\beta$ lies inside a non-trivial connection.
\end{lemma}

\begin{proof}

If $m = 0$, then necessarily $\beta = \sfrac{1}{2}$ and we have $c_{1/2} = \partial I_\alpha$ with $\pi_1(\alpha) \neq 1$. By \eqref{eq:sym_identity_endpoints_2}, $T^{-1}(c_{1/2}) = \partial I_{\hat \alpha}$, where $\pi_0( \hat \alpha) = \pi_0(\alpha) - 1$.

Without loss of generality, let us assume $m < 0$, the case $m > 0$ being analogous. By \eqref{eq:inverse_iterates},
\begin{equation*}
\label{eq:inverse_iterates_negative}
T^m(c_{\beta})=T^{-1}\circ \I_I(T^{-m-\delta_{1/2}(\beta)})\iff T^{-m- \delta_{1/2}(\beta) }(c_\beta) = \I_I\circ T \big( T^{m}(c_\beta)\big).
\end{equation*}
Notice that since $T^{m + 1}(c_\beta) \notin \{\partial I_\gamma\}_{\gamma \in \A}$ then $T^{m}(c_\beta) = \partial I_\alpha \neq \partial I_{\pi_0^{-1}(1)}$. Hence, by~\eqref{eq:sym_identity_endpoints}, the equation above implies $T^{-m - \delta_{1/2}(\beta)}(c_\beta) = \partial I_{\hat \alpha}$, where \hbox{$\pi_0( \hat \alpha) = \pi_0(\alpha) + 1$}.
\end{proof}

The previous lemma immediately implies the following.

\begin{corollary}
\label{cor:centers_connections}
Let $T: I \to I$ be a symmetric IET. Then any non-trivial irreducible connection of $T$ contains at most one point from the set $\left\{c_\alpha \mid \alpha \in \A \cup \{\sfrac{1}{2}\} \right\}$. In particular, there exists $\alpha \in \A$ such that $c_\alpha$ does not belong to any connection.
\end{corollary}

The following result provides a `recipe' to construct symmetric intervals dynamically by using iterates of the endpoints of the exchanged intervals.

\begin{lemma}\label{lem: sym_endpoints}
{Let $T : I \to I$ be a symmetric IET.}
Let $\alpha\in\mathcal A\setminus\{\pi_0^{-1}(1)\}$ and $m< M(\alpha)$. Then
\begin{equation}\label{eq: connections}
\I_I\circ T(T^{-m}(\partial I_{\alpha}))={\I_I\circ (T^{-m+1}(\partial I_{\alpha}))}=T^m(\partial I_{\hat\alpha}),
\end{equation}
where $\pi_0(\hat\alpha)=\pi_0(\alpha)-1$.

In particular, the left-closed right-open interval with endpoints $T^{-m}(\partial I_{\alpha})$ and $T^m(\partial I_{\hat\alpha})$ is {$\beta$-symmetric for some $\beta\in\A$, while the left-closed right-open interval with endpoints $T^{-m+1}(\partial I_{\alpha})$ and $T^m(\partial I_{\hat\alpha})$ is $\sfrac{1}{2}$-symmetric}.

Moreover, $M(\alpha)= N(\hat\alpha)$.
\end{lemma}

\begin{proof}
First, notice that \eqref{eq: connections} is equivalent to \eqref{eq:sym_identity_endpoints} if $m = 1$. This implies that $\I_I(\partial I_\alpha) = T(\partial I_{\hat\alpha})$.
In the following, we~assume $m > 1$. If $m < M(\alpha)$, by~\eqref{eq:generalized_conjugacy},
\[
\I\circ T(T^{-m}(\partial I_{\alpha}))= \I \circ T^{-m + 1}(\partial I_\alpha) = T^{m - 1} \circ \I (\partial I_\alpha) = T^m(\partial I_{\hat \alpha}).
\]
This proves \eqref{eq: connections}, which easily implies that the symmetry properties of the intervals in the statement.

Assume now that $m=M(\alpha)$ and $T^{-m}(\partial I_{\alpha})=\partial I_{\beta}$ for some $\beta\in \A$. Since $M(\alpha)$ is the minimal number that makes a connection, by~\eqref{eq: connections} we have
\[
\I\circ T(T(\partial I_{\beta}))=\I\circ T(T^{-{m+1}}\partial I_{\alpha})=T^{m-1}(\partial I_{\hat\alpha}).
\]
Moreover, noticing that $T(\partial I_\beta) \notin \{\partial I_\alpha\}_{\alpha \in \A}$ since $m > 1$, by~\eqref{eq: standardconjugacy}
\[
\I\circ T(T(\partial I_{\beta}))=T^{-1}\circ\I\circ T(\partial I_{\beta})
=T^{-1}(\partial I_{\hat\beta}),
\]
where $\pi_0(\hat\beta)=\pi_0(\beta)+1$. Thus we get $T^{m}(\partial I_{\hat\alpha})=\partial I_{\hat\beta}$. This proves $N(\hat\alpha)\le M(\alpha)$. The opposite inequality is proved analogously, {by exchanging the roles of $T$ and $T^{-1}$.}
\end{proof}

Using the previous lemma, it~is not difficult to construct symmetric intervals as in the statement of Proposition \ref{prop:induced_symmetric}.

\begin{lemma}
\label{lem:sym_int_disjoint}
Let $T: I \to I$ be an ergodic symmetric IET. Then there exists $\beta \in \A$ such that $c_\beta$ is not part of any non-trivial connection from $T$ and, for any $\epsilon > 0$, there exists a $\beta$-symmetric subinterval $J\subseteq I$ disjoint from the connections of $T$ satisfying $|J| < \epsilon$.
\end{lemma}

\begin{proof}
By Corollary \ref{cor:centers_connections}, there exists $\beta \in \A$ such that $c_\beta$ is not part of any connection of $T$ and, by~Corollary \ref{cor: existenceofinfiniteorbit}, there exists $\alpha \in \A \setminus \{\pi_0^{-1}(1)\}$ such that $M(\alpha) = +\infty$.

Since $T$ is ergodic and hence minimal, there exists $m \geq 1$ such that
\[
\big|T^{-m}(\partial I_\alpha) - c_\beta \big| < \sfrac{\epsilon}{2}.
\]
By taking $\epsilon$ smaller if necessary, we~may assume that $\epsilon < \min_{\delta \in \A} |I_\delta|$ and that $(c_\beta - \epsilon, c_\beta + \epsilon)$ does not contain any point from any connection.

Then, by~Lemma \ref{lem: sym_endpoints}, the left-closed right-open subinterval $J$ with endpoints $T^{-m}(\partial I_\alpha)$, $T^{m}(\partial I_{\hat \alpha})$ is $\beta$-symmetric, where $\pi_0(\hat \alpha) = \pi_0(\alpha) - 1$. In particular $J \subseteq (c_\beta - \sfrac{\epsilon}{2}, c_\beta + \sfrac{\epsilon}{2})$.
\end{proof}

We are now in a position to prove Propositions \ref{prop:induced_symmetric} and \ref{prop:induced_centers}.

\begin{proof}[Proof of Proposition \ref{prop:induced_symmetric}]
Let $T = (\pi, \lambda), J, m, \alpha, \hat \alpha$ as in the statement of the proposition, with $J$ not necessarily disjoint from the connections of $T$.

By Lemma \ref{lem: sym_endpoints}, there exists $\beta \in \A \cup \{\sfrac{1}{2}\}$ such that $J$ is $\beta$-symmetric, where $\beta = \sfrac{1}{2}$ only if the endpoints of $J$ are of the form $T^{-m + 1}(\partial I_\alpha)$, $T^m(\partial I_{\hat \alpha})$. Then, by~Lemma \ref{lem:symmetric_condition}, the induced IET $T_J$ satisfies \eqref{eq: standardconjugacy}.

From now on, let us assume that $J$ does not contain any point from any connection.
By Lemma \ref{lem: alphabetafetrconnections}, it~follows that $T_J = (\pi^J, \lambda^J)$ is an IET on $d - d'$ intervals, where~$d'$ is the number of non-trivial irreducible connections of $T$. In view of Proposition~\ref{prop: unwinding}, there exists $\tilde\lambda\in\Lambda^{\mathcal A,I}_T$ (see \eqref{eq:nbhd_connections}) such that $\tilde{T}:=(\pi,\tilde\lambda)$ is symmetric, the IET $\tilde{T}_{\tilde J} = (\tilde \pi^J, \tilde \lambda^J)$ obtained by inducing $\tilde{T}$ to the interval $\tilde J$ with endpoints $T^{-m}(\partial \tilde I_{\alpha})$ and $T^{m}(\partial \tilde I_{\hat\alpha})$ has the same permutation as $T_J$, and $\tilde\lambda^J$ has intervals of rationally independent lengths.

Since $\tilde J$ is a symmetric interval for $\tilde T$, by~Lemma \ref{lem:symmetric_condition} the induced IET $\tilde T_J$ satisfies~\eqref{eq: standardconjugacy}. Thus, by~Corollary \ref{cor: symmetricintcont}, $\tilde{T}_{\tilde J}$ exchanges its maximal continuity intervals symmetrically, and therefore so does $T_J$.

Hence, to prove that $T_J$ is a symmetric IET, it~suffices to show that it possesses exactly $d - d'$ maximal continuity intervals. By replacing $T_J$ with $\tilde{T}_{\tilde J}$ if necessary, we~may assume that the intervals exchanged by $T_{J}$ are of rationally independent lengths. Note that since $\tilde\lambda\in\Lambda^{\mathcal A,I}_T$, this change would not affect the connection pattern. Then also the intervals of continuity of $T_J$ are of rationally independent lengths. Let $\{I^J_{\alpha}\}_{\alpha\in\mathcal \A_J}$ be the intervals exchanged by $T_J$ and let $\{\hat I^J_{\alpha}\}_{\alpha\in\mathcal C}$ be the maximal continuity intervals of $T_J$. Then, to finish the proof, it~suffices to show that $\#\mathcal C = \#\mathcal \A_J$.

For every $\gamma\in\mathcal C$ consider the point $\hat c_\gamma^J$, the center-point of the interval $\hat I^J_{\gamma}$. Since~$T_J$ interchanges the intervals of continuity symmetrically, we~have $T_J(\hat c_\gamma^J)=\I_J(\hat c_\gamma^J)$ and no other point satisfies this equation. Moreover, since the lengths of intervals exchanged by $T_J$ are rationally independent, we~also have $\hat c_\gamma^J \notin\{\partial I^J_{\alpha}\}_{\alpha\in\A_J}$.

\begin{claim}
\label{cl:backward_centers}
Let $\sigma \in \A \cup \{\sfrac{1}{2}\}$ with $\sigma \neq \beta$ and $\ell := b_J(c_\sigma) \geq 1$ be the first backwards return time of $c_\sigma$ to $J$, that is, such that $p_J(c_\sigma) = T^{-\ell}(c_\sigma)$. Then, the following dichotomy holds.
\begin{itemize}
\item either $\{c_\sigma, T(c_\sigma), \ldots, T^{\ell - \delta_{1/2}(\sigma)}(c_\sigma)\} \cap \{\partial I_\delta\}_{\delta \in \A} = \emptyset$ and
\[
p_J(c_\sigma) = \hat c_\gamma^J,\qquad T_J(\hat c_\gamma^J) = T^{\ell - \delta_{1/2}(\sigma)}(c_\sigma) = T^{2\ell - \delta_{1/2}(\sigma)}(\hat c_\gamma^J), \quad \text{for some } \gamma \in \mathcal C,
\]
\item or $\{c_\sigma, T(c_\sigma), \ldots, T^{\ell - \delta_{1/2}(\sigma))}(c_\sigma)\} \cap \{\partial I_\delta\}_{\delta \in \A} \neq \emptyset$ and
\[
p_J(c_\sigma) = \partial I_\gamma^J,\quad\text{for some } \gamma \in \A_J,
\]
\end{itemize}
where $\delta_{1/2}$ is given by \eqref{eq:delta_1/2}. Moreover, in the latter case, $c_\sigma$ lies in a non-trivial connection of $T$.
\end{claim}

\begin{proof}
First, let us assume that $\{c_\sigma, T(c_\sigma), \ldots, T^{\ell - \delta_{1/2}(\sigma)}(c_\sigma)\} \cap \{\partial I_\delta\}_{\delta \in \A} = \emptyset$.

By \eqref{eq:inverse_iterates} and Lemma \ref{lem: local_reflection},
$T^{-k}(c_\sigma) \in J$ if and only if $T^{k - \delta_{1/2}}(c_\sigma) \in J$,
for any $k$ such that $0 \leq k \leq \ell$. Hence, the first visit time of $c_\sigma$ to $J$
via $T$ is $\ell - \delta_{1/2}(\sigma)$, which implies
\[
T_J(p_J(c_\sigma)) = T^{\ell - \delta_{1/2}(\sigma)}(c_\sigma).
\]
Moreover, $p_J(c_\sigma) = T^{-\ell}(c_\sigma)$ is a fixed point of $\I_J \circ T_J$. Indeed, by~\eqref{eq:generalized_conjugacy} and Lemma~\ref{lem: local_reflection},
\begin{align*}
\I_J \circ T_J(p_J(c_\sigma)) &= \I \circ T \big( T^{\ell -
\delta_{1/2}(\sigma)}(c_\sigma) \big) \\
&= T^{-\ell + \delta_{1/2}(\sigma)}\circ
T^{-1} \circ \I (c_\sigma) = T^{-\ell + \delta_{1/2}(\sigma)} (c_\sigma) =
p_J(c_\sigma).
\end{align*}
Therefore, since $p_J(c_\sigma)$ is a fixed point of $\I_J \circ T_J$ and $T_J$ exchanges its continuity intervals symmetrically, $p_J(c_\sigma) = \hat c_\gamma^J$, for some $\gamma \in \mathcal C$.

Now, let us assume that $\{c_\sigma, T(c_\sigma), \ldots, T^{\ell - \delta_{1/2}(\sigma)}(c_\sigma)\} \cap \{\partial I_\delta\}_{\delta \in \A} \neq \emptyset$.

Let $k$ with $0 \leq k \leq \ell - \delta_{1/2}(\sigma)$ be the minimum such that $T^k(c_\sigma) = \partial I _\delta$ for some $\delta \in \A$. By Lemma \ref{lem:symmetric_connection}, $T^{-k - \delta_{1/2}(\sigma)}(c_\sigma) = \partial I_{\hat \delta}$, where $\pi_0(\hat \delta) = \pi_0(\delta) - 1$. Since $-k - \delta_{1/2}(\sigma) \geq -\ell$, it~follows that $p_J(\partial I_{\hat \delta}) = p_J(c_\sigma)$ which, by~definition of $T_J$ (see Section \ref{sc:induced_IETs} and Lemma \ref{lem: alphabetafetrconnections}), coincides with $\partial I_\gamma^J$, for some $\gamma \in \A_J$. Notice that, in~this case, $c_\sigma$ lies in a non-trivial connection of $T$.
\end{proof}

The claim above shows that $p_J$ maps the set $\{c _\sigma \mid \sigma \in \mathcal B\}$, where
\begin{multline*}
\mathcal{B} := \{\sigma \in \A \cup \{\sfrac{1}{2}\} \mid \sigma \neq \beta \text{ and } c_\sigma \text{ does not belong to}\\
\text{any non-trivial irreducible connection}\},
\end{multline*}
injectively to the set $\{\hat c_\gamma^J \mid \gamma \in \mathcal C\}$.

Indeed, if for $\sigma, \sigma' \in \mathcal B$ we have $p_J(c_\sigma) = \hat c_\gamma^J = p_J(c_\sigma')$ then it follows from the previous claim that $r_J(\hat c_\gamma^J) = 2b_J(c_\sigma) - \delta_{1/2}(\sigma) = 2b_J(c_\sigma') - \delta_{1/2}(\sigma')$. Hence, since all the terms in the previous equality must have the same parity, it~follows that either $\sigma = \sfrac{1}{2} = \sigma'$ or $\sigma, \sigma' \in \A$. In the latter case, it~follows from the claim that $c_\sigma = T^{r_J( \hat c_\gamma^J)/2}(\hat c_\gamma^J) = c_\sigma'$.

Therefore, Claim \ref{cl:backward_centers} implies that $\# \mathcal C \geq \# \mathcal B$. Since, by~Corollary \ref{cor:centers_connections}, any non-trivial irreducible connection contains at most one point of the form $\{c_\sigma \mid \sigma \in \A \cup \{\sfrac{1}{2}\}\}$ it follows that $\#\mathcal B \geq d - d'$. Therefore $\#\mathcal C \geq d - d'$, which together with $\#\mathcal C\le\#\mathcal \A_J = d - d' \le\#\mathcal A$, implies
\[
\# \mathcal C = \#\mathcal B = d - d' = \# \A_J.\qedhere
\]
\end{proof}

\begin{proof}[Proof of Proposition \ref{prop:induced_centers}]
By Proposition \ref{prop:induced_symmetric}, it~follows that the maximal continuity intervals of $T_J$, which in the proof of Proposition \ref{prop:induced_symmetric} we denoted by $\{\hat I^J_{\alpha}\}_{\alpha\in\mathcal C}$, coincide with the intervals exchanged by $T_J$, which we denoted by $\{I^J_{\alpha}\}_{\alpha\in\mathcal \A_J}$.

Therefore, Proposition \ref{prop:induced_centers} follows from Claim \ref{cl:backward_centers} and the fact that $p_J$ maps the set
\begin{multline*}
\{c _\sigma \mid \sigma \in \A \cup \{\sfrac{1}{2}\},\, \sigma \neq \beta \text{ and } c_\sigma \text{ does not belong to}\\
\text{any non-trivial irreducible connection}\},
\end{multline*}
injectively to the set $\{\hat c_\gamma^J \mid \gamma \in \mathcal C\} = \{c_\gamma^J \mid \gamma \in \mathcal \A_J\}$, which we showed at the end of the proof of Proposition \ref{prop:induced_symmetric}.
\end{proof}

The following is a direct consequence of Proposition \ref{prop:induced_centers}.

\begin{corollary}\label{cor: emptyconnection}
If a symmetric IET $T$ has a non-trivial connection that does not contain a point from the set $\{c_{\alpha} \mid \alpha \in \A \cup \sfrac{1}{2}\}$, then $T$ is not ergodic.
\end{corollary}
{

\begin{proof} By Proposition \ref{prop:induced_centers}, we~have that the sets
\begin{multline*}
\{c _\sigma \mid \sigma \in \A \cup \{\sfrac{1}{2}\},\, \sigma \neq \beta \text{ and } c_\sigma \text{ does not belong to}\\
\text{any non-trivial irreducible connection}\},
\end{multline*}
and $\A_J$
have the same cardinality (see the end of the proof of Proposition \ref{prop:induced_centers}). Recalling that $d-\#\A_J$ is equal to the number of non-trivial irreducible connections and that each non-trivial irreducible connection contains at most one point from the set $\{c _\sigma \mid \sigma \in \A \cup \{\sfrac{1}{2}\}\} $ (see Corollary \ref{cor:centers_connections}), the Corollary follows by a simple counting argument.
\end{proof}}
\section{The essential values criterion}
In this section, we~recall and state the standard notion of essential value, which is a classical tool to study skew products' ergodicity. Let $(X,\mathcal B, \mu)$ be a standard probability space. Let $T:X\to X$ be $\mu$-measure preserving automorphism and let $f:X\to \R^m$, where $m \geq 1$. Consider the skew product $T_f$ on $X\times \R^n$ given by
\[
T_f(x,r)=(Tx,r+f(x)).
\]
We say that $a\in\R^m$ is an \emph{essential value} of $T_f$ if for every $\epsilon>0$ and every measurable subset $E\subseteq X$ with $\mu(E)>0$, there exists $n\in\N$ such that
\[
\mu\{x\in E\mid T^{n}(x)\in E\text{ and }|S_n f(x)-a|<\epsilon\}>0,
\]
where $S_n f (x)$ denotes the $n$-th Birkhoff sum of $f$ evaluated at $x$, which is given by
\begin{equation}
\label{eq:Birkhoff_sums}
S_n f (x):=\begin{cases}
\sum_{i=0}^{n-1}f(T^{i}(x))& \text{ if }n\ge 1,\\
0&\text{ if }n=0,\\
-\sum_{i=-n}^{-1}f(T^{i}(x))&\text{ if }n\le 1.
\end{cases}
\end{equation}
We denote the set of essential values of $T_f$ by $\Ess(T_f)$.

The following classical fact links this notion to the ergodicity of the skew product.
\begin{theorem}\label{thm: essvalues}
With the notation as above, the set $\Ess(T_f)$ is a closed subgroup of~$\R^m$. Moreover, $T_f$ is ergodic with respect to $\mu\otimes \Leb_{\R^m}$ if and only if $T$ is ergodic and $\Ess(T_f)=\R^m$.
\end{theorem}
We will use a simplified version of this criterion, as introduced by Conze and Frączek in
\cite{conze_cocycles_2011}.

\begin{theorem}[{\cite[Lem.\,2.7]{conze_cocycles_2011}}]\label{thm: FUcriterion}
Let $a\in\R^m$. Assume that for every $\epsilon>0$ there exists a sequence
of subsets $\{\Xi_n\}_{n\in\N}$ and an increasing sequence
$\{q_n\}_{n\in\N}$ of natural numbers such that
\begin{enumerate}
\item \label{cond:measure} $\liminf_{n\to\infty}\mu(\Xi_n)>0$,
\item \label{cond:rigidity}$\lim_{n\to\infty}\sup_{x\in \Xi_n}|T^{q_n}(x)-x|=0$,
\item \label{cond:almost_invariant} $\lim_{n\to\infty}\mu(\Xi_n\triangle T(\Xi_n))=0$,
\item \label{cond:BS} $|S_{q_n}f(x)-a|<\epsilon$.
\end{enumerate}
Then $a\in \Ess(T_f)$.
\end{theorem}

Lemma 2.7 in \cite{conze_cocycles_2011} actually states that
the
topological support of the limit distribution
$P:=\lim_{n\to\infty}\frac{1}{\mu(\Xi_n)}(S_{q_n}f(x)|_{\Xi_n})_*\mu|_{\Xi_n}$,
which exists up
to taking a subsequence due to tightness guaranteed by Condition
\eqref{cond:BS}, is contained in $\Ess(T_f)$. However, again by \eqref{cond:BS},
the topological support of $P$ is contained in $\overline{B(a,\epsilon)}\subset \R^m$. By
passing with $\epsilon$ to 0 and by the fact that $\Ess(T_f)$ is a closed subset
of $\R^m$, we~get that $a\in \Ess(T_f)$.

We now provide a version of the above criterion, that is going to be effective for our purposes.

\begin{proposition}\label{prop: effcriterion}
Let $T:I\to I$ be an ergodic IET and let $m\in \N_+$. Assume that the function $f:I^{\times m}\to \R^{m}$ satisfies the following:
\begin{enumeratei}
\item \label{cond:continuity} $f$ is of the form $(x_1,\ldots,x_m)\mto (f_1(x_1),\ldots,f_m(x_m))$ with each $f_j$ being differentiable over the exchanged intervals,
\item \label{cond:sequences} there exist sequences $\{\Xi_n\}_{n\in\N}$ and $\{q_n\}_{n\in\N}$ satisfying \eqref{cond:measure}-\eqref{cond:almost_invariant} in Theorem \ref{thm: FUcriterion}, with the additional assumption that, for any $n \in \N$, $\Xi_n=\bigsqcup_{i=0}^{h_n-1}T^{i}({I_n})$ is a Rokhlin tower of $h_n\le q_n$ intervals with $|I_n|\le\sfrac{1}{q_n}$ and $\sup_{x\in \Xi_n}|T^{q_n}(x)-x|\le D/q_n$ for some $D>1$,
\item \label{cond:cont_interval} for every $x\in\Xi_n$, the interval $[x,T^{q_n}(x)]$ (or $[T^{q_n}(x),x]$) is a continuity interval of $f_j$, for every $j\in\{1,\ldots,m\}$,
\item \label{cond:derivative} there exists $C>1$ such that
\[
|S_{q_n}f_j(x)|\le C, \qquad C^{-1}q_n\le S_{q_n}f_j'(x)\le Cq_n,
\]
for any $j=1,\ldots,m$ and any $x\in\Xi_n$, and
\[
|f_j'(x)|\le C,
\]
for any $x\in I$.
\end{enumeratei}
If additionally $T^{\otimes m}$ is ergodic, then so is $T_f^{\otimes m}$.
\end{proposition}

\begin{proof}
In view of Theorem \ref{thm: essvalues}, it~is enough to show that we can
obtain an $m$\nobreakdash-dimensional cube of essential values. We will prove the
proposition for $m=1$ by realizing the essential values through Rokhlin
towers, which are subsets of $\Xi_n$. For $m\ge 2$ we first obtain the
same result for every $j=1,\ldots,m$, \ie we find an interval $(a_j,b_j)$
of essential values, realized through Rokhlin towers inside $\Xi_n$. Hence,
the set $(a_1,b_1)\times\ldots\times(a_m,b_m)$ is the set of essential
values realized via sequences of Rokhlin towers inside
$\bigsqcup_{i=0}^{h_n-1}(T^{\otimes m})^i(I_n^{\times m})$.

That being said, we~assume from now on that $m=1$. Note that due to our assumption on $\Xi_n$ and the derivative of $f$, each of the sets $S_{q_n}f(\Xi_n)\subseteq [-2C,2C]$ is a uniformly bounded union of intervals. Note that, since
\[
\liminf_{n\to\infty} h_n |I_n| = \liminf_{n\to\infty}\mu(\Xi_n) > 0
\]
and $h_n\le q_n$, we~have that, by~passing to a subsequence if necessary, there exists $E>1$ with
\[
E^{-1}\le q_n|I_n|\le 1.
\]
Hence, given the assumption on the derivative, we~have
\[
\Leb_{\R}\left(S_{q_n}f(\Xi_n)\right)\ge \frac{1}{CE}>0
\]
for every $n\in\N$.
By passing to a subsequence if necessary, we~may assume that the sets $S_{q_n}f(\Xi_n)$ have a nonempty intersection. Let $y\in \bigcap_{i=1}^{\infty} S_{q_n}f(\Xi_n)$. We claim that~$y$ is an essential value. For this purpose, we~will construct a subtower $\Xi_{n}^{y}\subseteq \Xi_n$, i.e., a~subset of $\Xi_n$ which itself is a Rokhlin tower of intervals, depending on $\epsilon>0$, such that \eqref{cond:measure} and \eqref{cond:BS} in Theorem \ref{thm: FUcriterion} is satisfied. This is enough, since \eqref{cond:rigidity} and \eqref{cond:almost_invariant} are automatically satisfied by any sequence of subtowers of $\Xi_n$.

Fix $\epsilon>0$ and consider the point {$x_n \in \Xi_n$ such that $S_{q_n}f(x_n)=y$}.
Let $\ell_n\in\{0,\ldots,h_n-1\}$ be such that $x_n\in T^{\ell_n}(I_n)$ and assume without loss of generality that $\ell_n<h_{n}/2$, the other case being treated symmetrically. Since each level of the tower is an interval then
\begin{align*}
\tag*{either}
\Bigl(x_n,x_n+\frac{1}{\max\{C, D,E\} q_n}\Bigr)\subseteq T^{\ell_n}(I_n),\\
\tag*{or} \Bigl(x_n-\frac{1}{\max\{C,D,E\} q_n},x_n\Bigr)\subseteq T^{\ell_n}(I_n).
\end{align*}
Again, without loss of generality, let us assume that it is the former. Consider the tower
\[
\Xi_n^y:=\bigsqcup_{i=0}^{\epsilon h_n/4CD} T^i\left(x_n,x_n+\frac{\epsilon}{2\max\{C,D,E\} q_n}\right)\subseteq \Xi_n.
\]
We now show that for every $x\in \Xi_n^y$ we have $S_{q_n}f(x)\in(y-\epsilon,y+\epsilon)$. If $x\in T^{\ell_n}(I_n)$, then by the mean value theorem, we~have
\[
\left|S_{q_n}f(x)-y\right|=\left|S_{q_n}f(x)-S_{q_n}f(x_n)\right|\le Cq_n|x-x_n|\le \epsilon/2.
\]
If $x\in T^{j} (T^{\ell_n}(I_n))$ with $j=1,\ldots,\epsilon h_n/4CD$, then we
split the Birkhoff sum into two pieces $S_{q_n}f(x) = S_{q_n - j}f (x) +
S_jf(T^{q_n - j}(x))$, which we estimate separately. For the first term, we~have
\[
\left|S_{q_n-j} f(x)-S_{q_n-j} f(T^j(x_n))\right|\le Cq_n|x-T^j(x_n)|\le\epsilon/2.
\]

For every $k\in\{0,\ldots,h_n/2\}$ we have that $T^{k}(x_n)$ and
$T^{k}(T^{-j}(x))$ belong to the same level of $\Xi_n$ and, since $\Xi_n$ is a
tower of intervals as such, belong to the same interval continuity of $f$.
Moreover, by~\eqref{cond:cont_interval}, $T^k(T^{q_n-j}(x))=T^k(T^{q_n}(T^{-j}(x)))$ and
$T^{k}(T^{-j}(x))$ belong to the same
continuity interval of $f$ for every $k\in\{0,\ldots,h_n/2\}$. Hence
$T^k(T^{q_n-j}(x))$ and $T^k(x_n)$ belong to the same continuity intervals of
$f$. Hence, by~\eqref{cond:BS} and mean value theorem we
get
\[
\left|S_j f(T^{q_n-j}(x))-S_j f(x_n)\right|\le \frac{C\epsilon h_n}{4D}\,|T^{q_n-j}(x)-x_n|\le \frac{C\epsilon h_n}{4CD} \frac{2D}{q_n}\le\epsilon/2
\]
and thus $|S_{q_n} f(x)-y|\le \epsilon$. It~remains to notice that
\begin{align*}
\liminf_{n\to\infty}\Leb(\Xi_n^y)&\ge \liminf_{n\to\infty}\frac{\epsilon h_n}{4CD}\frac{1}{\max\{C,D,E\}}\,|I_n|\\
&\ge \liminf_{n\to\infty} \frac{1}{4(\max\{C,D,E\})^2}\Leb(\Xi_n)>0.
\end{align*}
Thus we have proved that $y\in \Ess(T_f)$. Note that for every $n\in\N$ {we have}
\[
x\in \Bigl(x_n,x_n+\frac{1}{2\max\{C,D,E\} q_n}\Bigr).
\]
Hence,
\[
S_{q_n}f(x)-S_{q_n}f(x_n)\ge C^{-1}q_n |x-x_n|.
\]
Therefore, by~applying similar reasoning as to $y$, we~obtain that every
\[
z\in\Bigl[y,y+\frac{1}{2C^{-1}\max\{C,D,E\}}\Bigr]
\]
is an essential value of $T_f$. This finishes the proof of the proposition.
\end{proof}
\section{Proofs of main results.}

This section contains the proof of Theorems \ref{thm: ergo}, \href{https://jep.centre-mersenne.org/item/10.5802/jep.302.pdf#page=4}{1.2} and \href{https://jep.centre-mersenne.org/item/10.5802/jep.302.pdf#page=5}{1.3}. In all three proofs, we~will apply the ergodicity criterion described in Proposition \ref{prop: effcriterion}. For this reason, we~start this section by outlining a construction that will be common to all of the proofs since it concerns only the underlying IET $T$ and not the cocycle being considered. At the end of this construction, we~will describe in detail why the assumptions of Proposition \ref{prop: effcriterion} are fulfilled in each setting.

Throughout this section, let $T=(\pi,\lambda): I \to I$ be an ergodic symmetric IET on $d = \#\A$ intervals $\{I_\alpha\}_{\alpha \in \A}$.

By Corollary \ref{cor:centers_connections}, there exists $\beta \in \A$ such that $c_{\beta}$ is not a part of any connection of~$T$. By Lemma \ref{lem:sym_int_disjoint}, there exists $\alpha \in \A$ and a nested sequence of $\beta$-symmetric intervals $\{J_n\}_{n \in \N}$ disjoint from the connections of $T$, with endpoints $T^{-m_n}(\partial I_{\alpha})$ and $T^{m_n}(\partial I_{\hat\alpha})$ for some $m_n \nearrow \infty$, where $\pi_0(\hat\alpha) = \pi_0(\alpha)-1$, and such that $\{c_\beta\} = \bigcap_{n \in \N} J_n$.

By Proposition \ref{prop:induced_symmetric}, for every $n\in\N$, the induced IET $T_{J_n}$ is a symmetric IET with $d - d'$ intervals, where $d'$ is the number of non-trivial irreducible connections of~$T$. Without loss of generality we may index their exchanged intervals using the same alphabet $\B \subseteq \A$. Let us denote by $\{I_\gamma^n\}_{\gamma \in \B}$ the intervals exchanged by $T_{J_n}$ and by $\{c_\gamma^n\}_{\gamma \in \B}$ their middle points, where, to avoid the use of double subscripts, we~changed our usual notation $I_\gamma^{J_n}$ (\resp $c_\gamma^{J_n}$) to $I_\gamma^n$ (\resp $c_\gamma^n$).

By Proposition \ref{prop:induced_centers}, each of the towers in the tower decomposition of $I$ associated with $T_{J_n}$
\begin{equation}
\label{eq:decomposition_sequence}
I = \bigsqcup_{\gamma \in \B} \bigsqcup_{i = 0}^{h_\gamma^n - 1} T^i(I_\gamma^n),
\end{equation}
where $h^n = (h^n_\gamma)_{\gamma \in \B}$ is some vector in $\N^\A_+$ (see Section \ref{sc:induced_IETs}), contains exactly one point from the set
\begin{multline*}
 \{c _\sigma \mid \sigma \in \A \cup \{\sfrac{1}{2}\},\, \sigma \neq \beta \text{ and } c_\sigma \text{ does not belong to}\\
\text{any non-trivial irreducible connection}\},
\end{multline*}
in the middle of its \emph{central level} $T^{\lfloor h_\gamma^n/2
\rfloor}(I^n_\gamma)$. More precisely, for every $\gamma \in \B$ there exists
$\sigma$ in the set above such that the middle point $c_\gamma^n$ of the
interval $I_\gamma^n$ verifies $c_\sigma = T^{\lfloor h_\gamma^n/2
\rfloor}(c_\gamma^n)$.

Using these facts, we~will show how to build towers $\{\Xi_n\}_{n\in\N}$ and a sequence $\{q_n\}_{n\in\N}$ that satisfies the assumptions of Proposition \ref{prop: effcriterion}.

A natural approach would be to consider subtowers of the already constructed towers. However in their current form, the towers may be very unbalanced: the wide towers may be very short and thus of very small measure, while thin towers may be very tall and contain the majority of the interval $I$. Since we need to construct towers of measure bounded away from 0, we~would have to choose them to be inside of the thin towers. This however makes it very difficult to control the rigidity of $T$ inside those towers as well as to estimate the values of Birkhoff sums. We tackle this problem by jumping between the points around which we induce.

Consider a Rokhlin tower $X_n:=\bigsqcup_{i=0}^{h_{\gamma_n}^n-1}T^i(I^n_{\gamma_n})$, where $\gamma_n \in\mathcal \B$ is chosen so that its Lebesgue measure is the largest compared to the other towers in the decomposition \eqref{eq:decomposition_sequence}. In particular
\[
\Leb(X_n)\geq \frac{1}{\#\B}\ge\frac{1}{\#\A}.
\]

Let us denote by $\mathfrak J_n$ the central level of this tower and recall that it contains a point $c_\sigma$ for some $\sigma \in \A \cup \{\sfrac{1}{2}\}$. By Proposition \ref{prop:induced_symmetric}, the induced transformation $T_{\mathfrak J_n}$ is a symmetric IET with $d - d'$ intervals, and we denote its exchanged intervals by $\{\mathfrak I_{\gamma}^n\}_{\gamma \in \mathcal B}$. As before, we~have a decomposition in Rokhlin towers of the form
\[
 I = \bigsqcup_{\gamma \in \B} \bigsqcup_{i = 0}^{\mathfrak h_\gamma^n - 1} T^i(\mathfrak I_\gamma^n).
\]

Let $\Gamma_n \in \B$ be such that $\mathfrak I_{\Gamma_n}^n$ is the largest of all intervals exchanged by $T_{\mathfrak J_n}$. Up to taking a subsequence, let us assume without loss of generality that there exists $\Gamma \in \B$ such that $\Gamma_n = \Gamma$, for every $n \in \N$.

As before, by~Proposition \ref{prop:induced_centers}, the tower $\bigsqcup_{i=0}^{\mathfrak h_{\Gamma}^n-1}T^i(\mathfrak I_{\Gamma}^n)$ contains exactly one point of the form $c_{\sigma}$ for some $\sigma \in \A \cup \{\sfrac{1}{2}\}$ in the middle of its central level. Let us denote this point by $\mathfrak c_n$.

Define
\[
\Xi_n:=\bigsqcup_{i=0}^{h_{\gamma_n}^n/2-1}T^{i}(\mathfrak I_{\Gamma}^n)\qquad \text{and}\qquad q_n:=\mathfrak h_{\Gamma}^n.
\]

Before passing to the proofs of each of the theorems, let us show that $\{\Xi_n\}_{n\in\N}$ and $\{q_n\}_{n\in\N}$ satisfy the assumptions \eqref{cond:sequences} and \eqref{cond:cont_interval} in Proposition \ref{prop: effcriterion}.

We start by showing that \eqref{cond:sequences} in Proposition \ref{prop: effcriterion} is satisfied, that is, that the sequences above verify \eqref{cond:measure}--\eqref{cond:almost_invariant} in Theorem \ref{thm: FUcriterion}.

First, we~can easily check that $\{\Xi_n\}_{n \in \N}$ satisfies \eqref{cond:measure} in Theorem \ref{thm: FUcriterion}. Indeed,
\begin{align*}
\Leb(\Xi_n) & =(h_{\gamma_n}^n/2)|\mathfrak I_{\Gamma}^n| \\ & > \frac{1}{\#\A}(h_{\gamma_n}^n/2)|\mathfrak J_n|=\frac{1}{\#\A}(h_{\gamma_n}^n/2)|I^n_{\gamma}| \\
& > \frac{1}{3\#\A} h_{\gamma_n}^n|I^n_{\gamma}|=\frac{1}{3\#\A} \Leb(X_n) \\
& >\frac{1}{3\#\A^2}.
\end{align*}

To see that $\{\Xi_n\}_{n\in\N}$ satisfies \eqref{cond:almost_invariant} in Theorem \ref{thm: FUcriterion} it is enough to notice that
\[
\mu(\Xi_n\triangle T(\Xi_n))\le 2|\mathfrak I_{\Gamma}^n|\to 0 \quad\text{as }n\to\infty.
\]

We will now show that
\begin{equation}\label{eq: derbound}
\sup_{x\in\Xi_n}|T^{q_n}(x)-x|<D/q_n\quad\text{for some }D>1
\end{equation}
thus showing \eqref{cond:rigidity} in Theorem \ref{thm: FUcriterion} as well. The argument uses the same observation as the one used to prove \eqref{cond:cont_interval} in Proposition \ref{prop: effcriterion}, which is the following.

Let $j\in\{0,\ldots,h_{\gamma_n}^n/2\}$ and let $x\in T^j(\mathfrak I_{\Gamma}^n)$. Then $T^{q_n-j}(x)\in \mathfrak J_n$. However, $\mathfrak J_n$ is the middle level of the tower $X_n$. Hence we have that $x$
and $T^{q_n}(x)=T^{j}(T^{q_n-j}(x))$ belong to the same level of the tower $X_n$. In particular, they belong to a continuity interval of $T$ (and as such to a continuity interval of any function continuous over exchanged intervals). Moreover, we~have
\[
|T^{q_n}(x)-x|\le |\mathfrak J_n|\le \#\A|\mathfrak I_{\Gamma}^n|=
\#\A\frac{1}{q_n} q_n|\mathfrak I_{\Gamma}^n|\le \frac{\#\A}{q_n}.
\]
Thus Conditions \eqref{cond:sequences} and \eqref{cond:cont_interval} in Proposition \ref{prop: effcriterion} are satisfied.

In the following proofs, we~will check that the remaining assumptions in Proposition~\ref{prop: effcriterion}, namely, Conditions \eqref{cond:continuity} and \eqref{cond:derivative}, are satisfied for the different cocycles considered in each setting. In view of the construction above this is enough to apply Proposition \ref{prop: effcriterion} and conclude the ergodicity of the skew product under consideration.

\begin{proof}[Proof of Theorem \ref{thm: ergo}]
We assume that $a>0$ since the other case follows symmetrically. We will apply Proposition \ref{prop: effcriterion} for $m=1$. Since $f(x)=a(x-\sfrac{1}{2})$ is continuous, \eqref{cond:continuity} in \ref{prop: effcriterion} is satisfied. We now show that \eqref{cond:derivative} is satisfied.

First, trivially, we~have
\[
f'(x)=a\quad\text{ for }x\in I.
\]
In particular
\[
S_{q_n} f'(x)=aq_n\quad \text{ for }x\in I.
\]
Recall that the tower $\bigsqcup_{i=0}^{\mathfrak h_{\Gamma}^n-1}T^i(\mathfrak I_{\Gamma}^n)$ has $\mathfrak c_n$ as its central point. By construction, the first visit time of $\mathfrak c_n$ via $T^{-1}$ to $\mathfrak J_n$ is $q_{n}/2$ or $(q_n+1)/2$. In both cases, by~Lemma \ref{lem: berktrujillo},
\[
S_{q_n}f(T^{-\lfloor (q_n+1)/2\rfloor}(\mathfrak c_n))=0.
\]
Since $S_{q_n}f$ is continuous in $\mathfrak I_{\Gamma}^n$ and $q_n|\mathfrak I_{\Gamma}^n|<1$, by~the mean value theorem we have that
\[
\left|S_{q_n}f(x)\right|<a\quad \text{ for every }x\in \mathfrak I_{\Gamma}^n.
\]
If $x\in T^j(\mathfrak I_{\Gamma}^n)$ for $j\in \{1,\ldots,h_{\gamma_n}^n/2\}$, then by \eqref{eq: derbound} and, again, by~the mean value theorem,
\[
\left|S_{q_n}f(x)-S_{q_n}f(T^{-j}(x)) \right|=
\left|S_j f(T^{q_n-j}(x))-S_j f(T^{-j}(x)) \right|
\le a D.
\]
Thus we get
\[
\left|S_{q_n}f(x)\right|<a(D+1) \quad \text{ for every }x\in\Xi_n.
\]
Since the bound does not depend on $n$, \eqref{cond:derivative} is satisfied and, by~Proposition \ref{prop: effcriterion}, the skew product $T_f$ is ergodic.
\end{proof}

\begin{thebibliography}{10}
\bibitem[2]{berk_ergodicity_2023} P. Berk, F. Trujillo \& C. Ulcigrai – “Ergodicity of explicit logarithmic cocycles over IETs”, 2023, arXiv:2210.16343.
\bibitem[3]{boissy_dynamics_2009} C. Boissy \& E. Lanneau – “Dynamics and geometry of the Rauzy–Veech induction for quadratic differentials”, Ergodic Theory Dynam. Systems 29 (2009), no. 3, p. 767–816.
\bibitem[4]{conze_cocycles_2011} J.-P. Conze \& K. Frączek – “Cocycles over interval exchange transformations and multivalued Hamiltonian flows”, Adv. Math. 226 (2011), no. 5, p. 4373–4428.
\bibitem[6]{veech_interval_1978} W. A. Veech – “Interval exchange transformations”, J. Analyse Math. 33 (1978), no. 1, p. 222–272.
\end{thebibliography}
\end{document}
